\documentclass[../main]{subfiles}
\ifSubfilesClassLoaded{\setcounter{chapter}{14}}{}
\begin{document}

\chapter{Evaluasi, Refleksi, dan Integrasi Kompetensi}

\begin{subcpmk}
  \item Sub-CPMK 13.1: Mengembangkan game lengkap dengan dokumentasi dan presentasi
  \item Sub-CPMK 13.2: Menerapkan best practices dan version control (Git)
\end{subcpmk}

Bab ini berisi asesmen komprehensif untuk mengukur pencapaian seluruh CPMK yang telah dipelajari sepanjang mata kuliah Pemrograman Game.

\section{Asesmen Akhir}

\subsection{Kuis Komprehensif}

Kuis mencakup semua bab: Konsep Game (MDA, GDD), Unity Engine, C\#, GameObject, Input, Physics, Animation, UI, Audio, Level Design, AI, Optimization, Build.

\subsection{Proyek Akhir}

Buat game lengkap (minimal setara Space Survivor) yang mendemonstrasikan:
\begin{itemize}
\item Player controller dengan input
\item Shooting dan collision
\item Score system dan UI
\item Audio (SFX, BGM)
\item Build untuk satu platform
\end{itemize}

\subsection{Presentasi}

Presentasikan proyek game: konsep, implementasi, tantangan, pembelajaran.

\section{Rubrik Penilaian}

Lihat Lampiran A untuk rubrik detail. Kriteria: Game Design, Implementasi Unity, Kode C\#, UI/UX, Audio, Build, Dokumentasi.

\section{Refleksi Pembelajaran}

\textbf{Pertanyaan Refleksi}:
\begin{enumerate}
\item Konsep mana yang paling menantang?
\item Bagaimana mengatasi kesulitan teknis?
\item Apa yang akan Anda pelajari lebih lanjut?
\item Bagaimana keterkaitan CPMK dengan proyek Anda?
\end{enumerate}

\section{Checklist Pencapaian Kompetensi Akhir}

\begin{checklist}
  \item Saya dapat menganalisis game menggunakan MDA framework
  \item Saya dapat mengoperasikan Unity Editor
  \item Saya dapat menulis skrip C\# untuk Unity
  \item Saya dapat mengimplementasikan gameplay (input, physics, collision)
  \item Saya dapat membuat UI dan Audio
  \item Saya dapat melakukan optimasi dan build
  \item Saya telah menyelesaikan proyek game
\end{checklist}

\begin{rangkuman}
Bab ini mengintegrasikan seluruh kompetensi melalui asesmen akhir, proyek, dan refleksi. Setelah menyelesaikan buku ajar ini, mahasiswa diharapkan mampu mengembangkan game menggunakan Unity Engine dan C\#. Terus praktik dan eksplorasi untuk mengembangkan keterampilan lebih lanjut.
\end{rangkuman}

\ifSubfilesClassLoaded{
  \renewcommand{\bibname}{Daftar Pustaka}
  \bibliographystyle{plain}
  \bibliography{../references}
}{}
\end{document}
